\section{How Infinity Could Emerge}

If the universe is infinite, or tends toward infinity, how could that come to be, and what would it even be like.
This chapter is the imaginative heart of the eternal story, drawn out carefully. The aim is to make infinity
thinkable, not to decide it. The question stays open. What follows is a set of ways infinity can enter a theory, a
crisp account of the cleanest one, an argument that infinity may be the more natural default, and a sober look at
what an infinite universe would feel like from inside.

\subsection{Seven ways infinity could emerge}

There is no single way for infinity to be. There are several, and they apply to different kinds of infinity, in
extent, in duration, in divisibility, and in cardinality. Most of them push infinity to the edges of the picture
rather than placing it in the middle.

\begin{itemize}
  \item As an unbounded process, that is, potential infinity. Infinity is never created. It is approached. A rule
  runs forever, finite at every step, and infinity is the name for the open horizon it tends toward. What is real
  is only the finite stages.
  \item As given and eternal, that is, actual infinity. Infinity is not created at all. It simply is, complete,
  with no origin. The cost is the classic paradoxes, Hilbert's Hotel and the rest.
  \item As a boundary, through compactification. Infinity is added to a finite-looking space as one tidy edge, the
  limit all outward paths share. You do not travel to it. This is literally the $S^3$ boundary at infinity of the
  bulk.
  \item As the continuum from the discrete, through coarse-graining. A discrete set has no points between its
  elements. Make it dense and complete it, and you get a continuum, infinitely many points in any interval. The
  base stays discrete and the continuum is the emergent large-scale description.
  \item As exponential or hyperbolic growth. Negative curvature opens space up so fast that it becomes effectively
  unbounded almost at once, the bridge from a tiny seed to an effectively infinite extent. This is exactly the
  exponential growth of the $\{3,4,3,4\}$ bulk.
  \item From a finite beginning, via creation and the arrow. A finite start, then unbounded growth driven by the
  arrow. The infinity is the future the process heads toward, not a thing that was made. This pairs the unbounded
  process with the arrow as the reason it runs forward.
  \item As the totality of possibilities, in ensembles. The infinite is not in any one world but in the collection
  of all of them.
\end{itemize}

The synthesis is that infinity is pushed to the edges. It is a limit of a process, a boundary object, or an
emergent description. It is rarely a completed actual thing sitting in the middle of reality.

\subsection{Finite but growing, the cleanest account}

The crispest version is this. At any single instant the universe has finitely many parts, but over all time the
count grows without bound. The universe is never actually infinite. It only tends toward infinity, and that
tending never finishes. Infinity is the open-ended horizon of the growth, not a place the universe arrives at.
Whether the completed limit is real or only a way of speaking is the metaphysical fork. The mathematics is
identical either way.

\subsection{Infinity as the default}

There is a stronger possibility, that infinity is the default. An eternal, beginningless, edgeless universe could
be the natural state, with no first moment to explain. The reversible rule makes this self-consistent. Started
from a generic state, it churns forever with steady activity and no decay, and it is exactly reversible, so every
state has a unique past and the history extends backward forever. No first moment is needed. On this view infinity
may be the more natural option for four reasons. An infinite homogeneous mesh has no privileged point, while a
finite one has a center or seed that demands explanation. An infinite eternal universe needs no boundary and no
special initial condition, while a finite one needs an edge or a starting state to account for. An infinite
eternal churn is already-everything, a natural attractor where the growing state is always incomplete. And
infinite available expanse can mean unbounded potential rather than a completed actual infinity, which slides back
toward finite-but-growing with no first moment.

This stays computable. Actual infinite extent is fine for a local rule as long as every finite region is finitely
computable. You never compute the whole at once. What is avoided is a completed infinity as a single graspable
object, and the continuum. An infinite past or an infinite mesh is allowed. An actual completed totality is not.

\subsection{The spectrum of positions}

Finite or infinite is too coarse a distinction. It is a spectrum, and a theory can be finite on one axis and
infinite on another, across extent, duration, divisibility, cardinality, and information. The positions form a
ladder.

\begin{table}[ht]
\centering
\begin{tabular}{@{}ll@{}}
\toprule
position & claim \\
\midrule
finitism & only finite objects exist, no actual infinity of any kind \\
finite and bounded & finite, closed, no edge, like a sphere's surface \\
finite but growing & finite at every moment, unbounded over time \\
potentially infinite & reality is an unbounded process, infinity its never-completed limit \\
actually infinite (extent) & space is infinite in size now \\
actually infinite (continuum) & the real continuum is a completed totality, uncountable \\
infinite ensembles & infinitely many worlds or all structures exist \\
\bottomrule
\end{tabular}
\caption{The spectrum from finite to infinite, across several axes.}
\end{table}

This substrate, the causal-set programs, Wolfram's project, and de Sitter cosmology all sit at the finite but
growing rung. The modern drift is downward and back, away from the view that infinity is actual and fundamental,
toward the view that infinity is potential, a boundary, or emergent. The base is finite and discrete, and the
continuum coarse-grains out of it.

\subsection{What it would feel like}

This is the part to sit with. A process can be unbounded and always finite at once. Counting is the model. There
is no largest number you will ever reach, yet every number you reach is finite. No largest does not mean an
infinite one exists. It means that for each, there is a bigger one. That is the felt sense of an endlessly growing
universe, always finite, never finished.

No edge does not mean infinite. Walk any direction on a sphere's surface and you come back, with no wall, no
boundary, and no outside, yet the surface is finite. A universe can be edgeless and still finite. Edgeless and
endless are different feelings, and keeping them apart removes most of the confusion.

A verb is not a noun. The claim that the universe is growing without end is potential infinity. The claim that the
universe contains infinitely many stars right now is actual infinity. The first feels like a becoming, the second
like a completed thing. The difference is completion, not size. Infinity itself is best pictured as a horizon
rather than a place. The boundary at infinity is a horizon of directions, a limit you approach and a screen you
read, not a faraway spot with room in it.

From inside, an eternal, beginningless, edgeless mesh would feel maximally symmetric and nothing-special. The rule
has simply always been running, every moment with a unique past behind it and a unique future ahead, a fire that
was never lit and never winds down. On the growing reading the texture is different. New moments are added at the
leading edge and stay. The now is the most recent slice of an ever-growing whole, time is lived duration, the
future is genuinely open, and novelty is real. The many become one and are increased by one.

\subsection{Where this substrate lands}

The working default is finite but growing, a hyperbolic crystal finite at every beat and unbounded over all
beats, discrete at the base with the continuum emergent. On that reading infinity is allowed only two
non-fundamental roles. It is the boundary the growth approaches but never reaches, and it is the emergent
continuum you see when you stop resolving cells. Everything real at every moment is finite. But the eternal,
beginningless, infinite picture is fully self-consistent too, since the reversible rule needs no first moment, so
it remains a live open axiom. The discriminator the analysis names is reversibility. A reversible rule needs no
beginning, which leaves the eternal reading allowed. An irreversible rule decays and needs a seed, which would
force a start. Our rule is reversible, so both doors stay open, exactly the fork left to the reader.
